\section{模型评价与结论}
本方法把逻辑生命与物理驻留区分开来，能够显式计算换出释放地址产生的同步关系；通过先生成再独立读回验证，得到覆盖三个问题的完整文本结果。固定驻留图重排保持换出集合不变，因而能够在严格零搬运增量下改善流水时间。这是本轮六图中最稳定的优化来源。

方法的主要不足在于候选空间有限。最大驻留目标的搜索尚不充分，五图未超过收缩后的编号基准；贪心分配无法证明搬运量最小，碎片恢复可能产生较多额外换出；固定物理复用顺序又限制了时间优化空间。后续可在保留完整约束重放的基础上，引入有界局部交换和更精细的连续区间牺牲代价，并在新增图集上评价泛化，而不依据本轮六图直接作普遍性能承诺。

六个实例的最终时间改善为 1.30\% 至 33.38\%，搬运量均未增加。该结论针对题面抽象执行模型，不等价于真实芯片上的测量加速。特别是换出期间地址复用采用本文的驻留段解释，尚未获得官方评测程序的交叉确认。

本文模型推导、程序和实验分析使用 OpenAI Codex 辅助完成，工具精确模型标识及版本发布日期未核实。所有数值来自本地运行与文件重放，尚未经参赛队员独立复核。正式参赛稿应依据适用的人工智能使用规定，由队员理解、核验并用自身语言修订\cite{ai}。
